\section{Introduction}

A guidance system makes two decisions, not one.

\begin{description}
\item[Level 1.] Which routing \emph{policy} should be active under the current
traffic conditions?
\item[Level 2.] Given that policy, which route does each guided vehicle receive?
\end{description}

Level 2 is the subject of a large literature. This paper concerns Level 1, and
treats the policies themselves as fixed, published objectives.

The question matters because established policies embody different assumptions.
Distance minimisation assumes the shortest corridor can absorb what is sent to
it. Reactive travel-time routing assumes recent measurements predict the near
future, which fails when a large guided cohort acts on them
simultaneously~\cite{arnott1991information,acemoglu2018braess}. Load balancing
assumes capacity is unevenly distributed and worth equalising. Reliability-aware
routing assumes travel-time variability is itself a cost. Each assumption holds
in some conditions and not others, so the policy that should be active
plausibly depends on the traffic --- and production systems nevertheless fix it
at design time.

\paragraph{What the literature establishes, and what it does not.}
Prior work supplies the routing objectives themselves, the traffic-feedback
mechanisms that make them fail, the algorithm-selection framework for choosing
among methods by instance features, the regret and hindsight-benchmark
reference points, and the uncertainty machinery for deciding when not to act.
Each is well developed, and Section~\ref{sec:lit} reviews them. What the
literature does not supply is a measurement of how much a context-dependent
choice among such policies is actually worth once they are placed in the same
physical environment and compared on the cost of the decision rather than on
the accuracy of a prediction.

\paragraph{The research gap.}
Showing that the preferred policy changes with conditions is not the same as
showing that a context-aware selector is worth building. A winner map can vary
substantially while the cost of ignoring it remains negligible, if the best
fixed policy loses only slightly wherever it is not best. These are separate
empirical quantities, and conflating them is easy because the first is the more
visible. A third question is separate again: whether a selector that captures
the available value can be trusted outside the conditions it was fitted on.
This paper measures all three.

\paragraph{Research questions.}
\textbf{RQ1} Under which traffic, signal, information, penetration, capacity and
disruption conditions do established routing policies exhibit complementary
strengths?
\textbf{RQ2} Are the preference boundaries explicable by mechanistic traffic
variables rather than only by a statistical model?
\textbf{RQ3} Does a learned selector identify useful selection structure beyond
a strong mechanistic rule?
\textbf{RQ4} When the operating context lies outside the training support, or
uncertainty is too large, can the system abstain safely?

\paragraph{Contributions.}
\begin{enumerate}
\item A factorial characterisation of the operating regions in which four
established routing policies are complementary, over \NumContexts{} contexts
and \NumRunsMain{} paired microsimulation runs.
\item A mechanistic account of the preference boundary in dimensionless traffic
descriptors, reported as brackets between sampled levels rather than as point
thresholds.
\item A decision-oriented evaluation in which the primary comparison is a
learned selector against a mechanistic rule, not against a fixed policy, and in
which selection accuracy and decision regret are shown to rank methods
differently.
\item A selective mechanism with a support gate and a confidence gate,
evaluated separately under interpolation, extrapolation, concept shift and
domain shift --- including one shift it does not detect.
\end{enumerate}

\paragraph{Scope and non-claims.}
We introduce no routing algorithm, no learning algorithm and no multi-criteria
decision method; every policy and model class used here is established and
cited. The environment is a controlled microsimulation, not a model of any
city, and no claim of real-world validation is made. We make no priority claim:
each component has precedent, and Section~\ref{sec:lit} states what is and is
not new.

\paragraph{What this paper demonstrates.}
In this benchmark the four policies are genuinely complementary on the
three-criterion outcome vector, yet the entire decision value of exploiting
that complementarity is \NumHeadroom~s of system mean journey time per vehicle
--- \NumHeadroomPct{} of the mean journey time achieved by the single best
fixed policy --- because that policy is near-dominant. Selecting the wrong
fixed policy costs \NumFixedPenPctAbsPOne{} of the same baseline. Within that
narrow margin, what a selector is trained to predict determines whether it
helps or hurts, and a selector that captures the available value is not
uniformly reliable once the operating conditions leave the training range.

\section{Literature Review}
\label{sec:lit}

\subsection{Dynamic and Traffic-Aware Routing}

Route choice under congestion rests on the equilibrium tradition founded by
Wardrop~\cite{wardrop1952}, in which drivers minimising their own travel time
reach a state that need not minimise total travel time. Shortest-path routing
on a static cost is the limiting case in which no traffic state is observed at
all; it is the zero-information reference against which any adaptive policy
must justify its infrastructure.

Reactive guidance --- routing on recently measured link travel times --- is the
operational descendant of that tradition, and its central difficulty is well
documented. Providing information to drivers need not reduce congestion,
because the response to the information changes the state it
described~\cite{arnott1991information}. Acemoglu et
al.~\cite{acemoglu2018braess} give this a sharp game-theoretic form, showing
that additional information can strictly worsen equilibrium outcomes. The
severity depends on how much of the demand responds, which makes guidance
penetration a first-order variable rather than an implementation
detail~\cite{yang1998penetration}, and on how stale the information is when it
is acted upon.

Two policy families respond to these failures directly. Load-balancing guidance
distributes vehicles over a set of travel-time-acceptable alternatives
according to predicted link load, rather than sending every vehicle to the
current best path; we use the flow-balanced $k$-shortest-path construction of
Pan et al.~\cite{pan2013proactive}. Reliability-aware guidance treats
travel-time variability as a cost in its own right, adding a dispersion term to
expected travel time, following the mean-variance formulation of Sen et
al.~\cite{sen2001meanvariance}. Both are established families with published
objectives, and we take them as given rather than modifying them.

\subsection{Environmental, Multi-Criteria, and Context-Dependent Routing}

Eco-routing minimises an emissions proxy rather than travel
time~\cite{boriboonsomsin2012eco}. Whether such an objective is \emph{distinct}
from a temporal one is a property of the network and the fleet, not of the
objective: where emissions are nearly monotone in travel time, an
environmental policy reproduces a time-based one under a different name. We
therefore treat the question as a measurement rather than an assumption, and
report the correlation and the frequency of disagreement rather than asserting
either.

More generally, comparing routing policies on a single scalar hides the
structure that multiple criteria expose. Where criteria disagree, the identity
of the preferred policy depends on which criterion is chosen, and a Pareto
analysis makes that dependence visible before any weighting is applied. We
adopt that ordering --- dominance first, scalarisation only as sensitivity ---
and avoid monetary valuation entirely, since no verified value of time, value
of reliability or carbon price was available for this setting; inventing one
would make the conclusion an artefact of the invented constant.

Context dependence enters when the operating conditions vary. Demand, signal
timing, guidance penetration, information lag, capacity distribution and
disruption each change which of the failure modes above is active. The
literature on route guidance establishes that these conditions matter for a
\emph{given} policy. What is less settled is whether they matter enough to
change \emph{which} policy should be running, and by how much.

\subsection{Algorithm Selection and Decision-Oriented Evaluation}

Choosing among algorithms by instance features is Rice's algorithm-selection
problem~\cite{rice1976}, with a mature methodology developed largely in
combinatorial search~\cite{xu2008satzilla,smithmiles2009,kotthoff2014survey}.
That literature supplies two reference points we adopt directly: the
\emph{single best solver}, the one fixed method with the lowest mean cost, and
the \emph{virtual best solver}, the retrospective per-instance
optimum~\cite{bischl2016aslib}. Instance-space analysis asks where in a feature
space each algorithm is strong~\cite{smithmiles2014instance}, which is
precisely the form RQ1 takes for routing policies.

Two further strands matter here. First, evaluating a predictor by the cost of
the decision it induces rather than by its predictive error is the
predict-then-optimise distinction of Elmachtoub and
Grigas~\cite{elmachtoub2022spo}; we show below that the two criteria rank
methods differently in this setting. Second, declining to decide is
classical~\cite{chow1970reject}, and split conformal prediction supplies
distribution-free intervals with which to decide when~\cite{vovk2005alrw,%
lei2018conformal}. Its coverage guarantee rests on exchangeability, which
covariate shift breaks~\cite{shimodaira2000,tibshirani2019covshift}. We treat
that as a design constraint rather than a caveat: a conformal gate is a
calibrated instrument in-distribution and a heuristic outside it.

\paragraph{Positioning.}
The components above exist separately and are each well developed: the routing
objectives, the feedback mechanisms that make them fail, the selection
framework, the regret and hindsight reference points, and the uncertainty
machinery. What is missing is a measurement that puts them together. This study
combines them into a single empirical policy-selection evaluation, run under a
controlled factorial of traffic conditions in a physical microsimulation, and
reports decision regret rather than selection accuracy as the primary quantity.

We do not claim that regime-dependent routing has not been studied, nor that no
prior work has combined any subset of these ideas. The contribution is the
measurement and its interpretation: how much decision value a context-dependent
choice among established policies carries when every component is held fixed and
known, and whether a selector capturing that value can be relied on outside the
conditions it was fitted on.
